\chapter{Introduction}
\label{ch:introduction}

Cybersecurity is no longer a niche technical concern. It directly affects daily life through
banking, communication, healthcare, transportation, and digital public services.
Modern attacks are targeted, persistent, and increasingly systemic because software systems
are deeply connected through cloud services, CI/CD pipelines, open-source dependencies,
and distributed teams.

This report presents a broad overview of contemporary cyber threats and then focuses on
software supply chain attacks as the central theme. The focus is motivated by recent incidents
where compromise of a trusted upstream component propagated silently to many downstream
consumers.

\section{Modern Cyber Threat Landscape}

The most common operational threat classes discussed in this report are malware/ransomware,
phishing and social engineering, denial-of-service (DDoS), insider threats, and software supply
chain attacks. A concise summary is provided in \cref{tab:threat-overview}.

\begin{table}[htbp]
\centering
\caption{High-level threat classes and representative mitigations.}
\label{tab:threat-overview}
\begin{tabular}{p{0.26\textwidth} p{0.30\textwidth} p{0.34\textwidth}}
\toprule
\textbf{Threat Class} & \textbf{Typical Attack Vector} & \textbf{Representative Mitigations} \\
\midrule
Malware and Ransomware & Malicious payload execution for encryption, theft, or disruption & Endpoint protection, offline backups, rapid patching \\
Phishing and Social Engineering & Deceptive communication that steals credentials or installs malware & MFA, secure email gateway, user awareness \\
DDoS & Flooding network/service resources to exhaust availability & Traffic filtering, rate limiting, redundancy \\
Insider Threats & Misuse of legitimate internal access & Least privilege, monitoring, segregation of duties \\
Supply Chain Attack & Compromise of trusted package, dependency, or build path & Provenance checks, signed releases, continuous vetting \\
\bottomrule
\end{tabular}
\end{table}

\section{Scope and Focus}

The broad threat overview is narrowed to supply chain attacks for three reasons:

\begin{enumerate}[label=(\roman*)]
\item the trust boundary in dependency ecosystems is often implicit,
\item one upstream compromise can impact many organizations,
\item detection and response are difficult when malicious behavior is hidden in normal update paths.
\end{enumerate}

\section{Key Contributions of This Report}

The main contributions are:

\begin{itemize}
\item a structured overview of modern cyber threats and mitigations;
\item a focused case study of the XZ Utils incident (2024) and its attack path;
\item synthesis of related work on package ecosystem security and dependency confusion;
\item in-depth discussion of existing and proposed defenses, with explicit limitations and future directions.
\end{itemize}

\section{Report Organization}

\Cref{ch:background} introduces background, core concepts, and related work.
\Cref{ch:casestudy} presents the XZ backdoor case and formal problem statement.
\Cref{ch:solutions} provides in-depth discussion of existing and proposed solutions.
\Cref{ch:conclusion} summarizes conclusions and outlines future work.


